% Generated by make_glossary.py from glossary.json. Do not edit by hand.
\newglossaryentry{acdf}{
  name={anterior cervical discectomy and fusion (ACDF)},
  description={Surgery through the front of the neck that removes a degenerated disc and any bone spurs behind it, then fuses the two vertebrae with a bone graft or cage and a plate. Suited to compression from the front at one or two levels.}
}
\newglossaryentry{adjacentsegment}{
  name={adjacent segment disease},
  description={New degeneration at the level next to a fusion, thought to be accelerated by the extra motion and load that the fused segment transfers to it. It can require further surgery years later.}
}
\newglossaryentry{als}{
  name={amyotrophic lateral sclerosis (ALS)},
  description={Motor neuron disease: progressive death of the nerve cells that drive muscles, causing weakness, wasting, and brisk reflexes without sensory loss. Its early hand weakness and mixed reflex pattern make it the most consequential mimic of cervical myelopathy.}
}
\newglossaryentry{autonomic}{
  name={autonomic dysfunction},
  description={Disturbance of the involuntary nervous control of bladder, bowel, and sexual function. In cervical myelopathy it usually appears as urgency, frequency, or hesitancy, and reflects damage to pathways running through the cord.}
}
\newglossaryentry{b12}{
  name={subacute combined degeneration},
  description={Damage to the dorsal columns and corticospinal tracts from vitamin B12 deficiency (or, less often, copper deficiency), producing numb hands and feet, unsteady gait, and brisk reflexes. It is treatable and is checked with a blood test before attributing symptoms to compression.}
}
\newglossaryentry{babinski}{
  name={Babinski sign},
  description={Stroking the sole of the foot makes the big toe extend upward instead of curling down. It is a very specific marker of damage to the motor pathways of the spinal cord or brain, but is absent in most people with cervical myelopathy.}
}
\newglossaryentry{bscb}{
  name={blood–spinal cord barrier},
  description={The tight lining of the cord's blood vessels that normally keeps blood proteins and immune cells out of nervous tissue. Compression makes it leaky, admitting inflammatory cells that add to the injury.}
}
\newglossaryentry{c5palsy}{
  name={C5 palsy},
  description={Weakness of the deltoid and biceps, the muscles supplied by the fifth cervical nerve root, appearing days after decompression surgery. It affects roughly 2\% of patients in large series and usually recovers over months.}
}
\newglossaryentry{carpaltunnel}{
  name={carpal tunnel syndrome},
  description={Compression of the median nerve at the wrist, causing numbness and tingling in the thumb, index, and middle fingers, often at night. It is the most frequent misdiagnosis in cervical myelopathy, which can produce similar hand symptoms in both hands at once.}
}
\newglossaryentry{clonus}{
  name={clonus},
  description={Rhythmic, involuntary muscle beating provoked by a sudden stretch, most often tested by briskly pushing the foot upward at the ankle. Like the Babinski sign it is highly specific for spinal cord injury but present in only about one patient in ten.}
}
\newglossaryentry{compressionratio}{
  name={compression ratio},
  description={The front-to-back thickness of the spinal cord divided by its side-to-side width on an axial MRI slice. A normal cord is roughly oval; a ratio of 0.4 or less marks a flattened cord and predicts progression to myelopathy.}
}
\newglossaryentry{cordcompression}{
  name={spinal cord compression},
  description={An MRI finding: something (disc, bone, or ligament) indents or flattens the spinal cord and displaces the fluid around it. Compression is necessary for degenerative cervical myelopathy but is not the same thing: about a quarter of healthy adults have it with no symptoms.}
}
\newglossaryentry{corpectomy}{
  name={corpectomy},
  description={Removal of most of a vertebral body, together with the discs above and below, through the front of the neck, followed by reconstruction with a strut and plate. Used when compression from the front extends behind the vertebral body itself.}
}
\newglossaryentry{corticospinaltract}{
  name={corticospinal tract},
  description={The main motor pathway running from the brain's motor cortex down the side of the spinal cord to the motor neurons that move the limbs. Damage to it below the neck produces weakness, spasticity, and brisk reflexes in the legs.}
}
\newglossaryentry{csm}{
  name={cervical spondylotic myelopathy (CSM)},
  description={The older, narrower name for spinal cord injury from spondylosis of the neck. It is now treated as the commonest subtype of degenerative cervical myelopathy, which also covers ligament ossification and congenital narrowing.}
}
\newglossaryentry{demyelination}{
  name={demyelination},
  description={Loss of the fatty insulating sheath (myelin) around nerve fibers, which slows or blocks conduction. In cervical myelopathy the fibers of the corticospinal tract and dorsal columns lose myelin at and around the compressed level.}
}
\newglossaryentry{dorsalcolumns}{
  name={dorsal columns},
  description={The bundles of nerve fibers at the back of the spinal cord that carry fine touch, vibration, and position sense upward to the brain. Their damage explains the clumsy hands and poor balance of cervical myelopathy.}
}
\newglossaryentry{dti}{
  name={diffusion tensor imaging},
  description={An MRI technique that measures how freely water moves along versus across nerve fibers. Reduced directional diffusion indicates fiber injury, and a 2025 AO Spine review conditionally recommends diffusion-weighted signal changes for judging prognosis after surgery.}
}
\newglossaryentry{dysphagia}{
  name={dysphagia},
  description={Difficulty swallowing. It is the commonest complication of front-of-neck operations, because the esophagus is moved aside to reach the spine, and is usually temporary.}
}
\newglossaryentry{emg}{
  name={electromyography (EMG)},
  description={Recording the electrical activity of muscles with a fine needle. It detects loss of nerve supply to a muscle and is mainly used to identify nerve-root involvement or to exclude motor neuron disease and peripheral neuropathy.}
}
\newglossaryentry{fusion}{
  name={spinal fusion},
  description={Joining two or more vertebrae so they heal into one bone, using bone graft held by screws, plates, or rods. It stabilizes the decompressed segment at the cost of motion and of extra load on neighboring levels.}
}
\newglossaryentry{grade}{
  name={GRADE recommendation strength},
  description={A system for grading clinical guidance. A strong recommendation means most patients should receive the intervention; a weak (conditional) recommendation means the choice depends on the patient's values and circumstances. Each is paired with a rating of the quality of the underlying evidence, from very low to high.}
}
\newglossaryentry{graymatter}{
  name={gray matter},
  description={The butterfly-shaped core of the spinal cord containing nerve cell bodies, including the motor neurons in its front horns that supply the arm and hand muscles at each level. Its loss at the compressed level causes wasting and weakness of the hands.}
}
\newglossaryentry{hazardratio}{
  name={hazard ratio (HR)},
  description={A relative risk that accounts for time: the instantaneous rate of an event in one group divided by the rate in another over the follow-up period. A hazard ratio of 1.57 means events occur about 57\% faster in the first group.}
}
\newglossaryentry{hoffmann}{
  name={Hoffmann sign},
  description={Flicking the tip of the middle finger downward makes the thumb and index finger flex involuntarily. Present in roughly 60\% of people with cervical myelopathy and in a smaller share of healthy people, so it raises suspicion without settling it.}
}
\newglossaryentry{hsp}{
  name={hereditary spastic paraplegia},
  description={A group of inherited disorders in which the long motor fibers to the legs slowly degenerate, producing progressive leg stiffness and gait difficulty. It mimics the leg findings of cervical myelopathy but has no cord compression and often a family history.}
}
\newglossaryentry{hyperreflexia}{
  name={hyperreflexia},
  description={Exaggerated tendon reflexes, such as a knee jerk that is brisker than normal. It is the most sensitive sign of cervical myelopathy but is also common in healthy people, so it supports rather than makes the diagnosis.}
}
\newglossaryentry{ibudilast}{
  name={ibudilast},
  description={An anti-inflammatory phosphodiesterase inhibitor being tested as an add-on to decompression surgery for degenerative cervical myelopathy in the UK RECEDE-Myelopathy trial.}
}
\newglossaryentry{icd}{
  name={International Classification of Diseases (ICD) code},
  description={The World Health Organization's coding system used in hospital records and statistics. Degenerative cervical myelopathy had no specific code, so cases are scattered across several codes and undercounted in administrative data.}
}
\newglossaryentry{intervertebraldisc}{
  name={intervertebral disc},
  description={The fibrous cushion between two vertebral bodies. With age it dries, flattens, and bulges backward toward the spinal canal, the first step in most cases of spondylosis.}
}
\newglossaryentry{invertedsupinator}{
  name={inverted supinator sign},
  description={Tapping the brachioradialis tendon near the wrist produces finger flexion instead of, or in addition to, the expected elbow flexion. It localizes cord compression to the C5 to C6 level and is highly specific.}
}
\newglossaryentry{ischemia}{
  name={ischemia},
  description={Shortage of blood supply to a tissue. Chronic compression pinches the small arteries that feed the spinal cord, and this starvation of oxygen is thought to drive much of the cell loss in cervical myelopathy.}
}
\newglossaryentry{jla}{
  name={James Lind Alliance priority setting partnership},
  description={A structured process in which patients, carers, and clinicians jointly rank the unanswered research questions that matter most to them. It was used to produce the ten research priorities for degenerative cervical myelopathy.}
}
\newglossaryentry{kyphosis}{
  name={kyphosis},
  description={Loss of the neck's normal curve so that the cervical spine bows the wrong way, convex toward the back. It drapes the cord over bone spurs at the front and is a reason to prefer front-of-neck surgery over a posterior approach.}
}
\newglossaryentry{laminectomy}{
  name={laminectomy},
  description={Removal of the laminae, the bony roof of the spinal canal, from the back of the neck to give the cord room. It is usually combined with screws and rods (fusion) to prevent the neck from later slumping forward.}
}
\newglossaryentry{laminoplasty}{
  name={laminoplasty},
  description={A posterior operation that hinges the laminae open like a door and props them in the widened position, enlarging the canal while keeping motion and avoiding fusion. Suited to compression at several levels in a neck that keeps its normal curve.}
}
\newglossaryentry{lhermitte}{
  name={Lhermitte sign},
  description={An electric-shock sensation running down the spine or into the limbs when the neck is flexed. It indicates irritation of the sensory pathways in the cervical cord and is seen in cervical myelopathy, multiple sclerosis, and vitamin B12 deficiency.}
}
\newglossaryentry{ligamentumflavum}{
  name={ligamentum flavum},
  description={The elastic yellow ligament lining the back wall of the spinal canal between adjacent vertebral arches. When it thickens with age or buckles as the neck extends, it presses on the cord from behind.}
}
\newglossaryentry{lmnsigns}{
  name={lower motor neuron signs},
  description={Findings of damage to the nerve cells or roots that directly supply muscle: wasting, weakness, reduced reflexes, and fasciculations. In cervical myelopathy they can appear in the arms at the level of compression, while the legs show upper motor neuron signs.}
}
\newglossaryentry{mcid}{
  name={minimum clinically important difference (MCID)},
  description={The smallest change in a score that patients notice as a real improvement. For the mJOA it is about 1 point in mild disease, 2 in moderate, and 3 in severe, which is why a small average gain in mild disease is hard to interpret.}
}
\newglossaryentry{mep}{
  name={motor evoked potentials (MEPs)},
  description={Muscle responses recorded after magnetic stimulation of the brain's motor cortex. A delayed response shows that the motor pathways in the spinal cord conduct slowly, and helps distinguish cord disease from motor neuron disease.}
}
\newglossaryentry{mjoa}{
  name={modified Japanese Orthopaedic Association score (mJOA)},
  description={An 18-point clinician-rated scale of arm function, leg function, hand sensation, and bladder control, where 18 is normal. It defines severity: mild is 15 to 17, moderate 12 to 14, severe 11 or less, and it is the main outcome measure in treatment trials.}
}
\newglossaryentry{mri}{
  name={magnetic resonance imaging (MRI)},
  description={The imaging method that shows soft tissue, including the spinal cord, discs, and ligaments, without radiation. It is the essential test for degenerative cervical myelopathy because X-rays and CT show bone but not the cord.}
}
\newglossaryentry{ms}{
  name={multiple sclerosis (MS)},
  description={An inflammatory disease in which the immune system strips myelin from patches of the brain, optic nerves, and spinal cord. A spinal cord plaque can copy cervical myelopathy, but MS typically starts younger, relapses, and shows brain lesions.}
}
\newglossaryentry{myelomalacia}{
  name={myelomalacia},
  description={Softening and loss of spinal cord tissue, the end result of long-standing compression and ischemia. It appears as a dark area on T1-weighted MRI and marks damage that surgery cannot reverse.}
}
\newglossaryentry{myelopathicsigns}{
  name={upper motor neuron signs},
  description={Examination findings that point to damage of the motor pathways in the brain or spinal cord rather than the peripheral nerves: brisk reflexes, spasticity, clonus, and pathological reflexes such as the Hoffmann and Babinski signs. In cervical myelopathy they appear below the level of compression.}
}
\newglossaryentry{myelopathy}{
  name={myelopathy},
  description={Any disease or injury of the spinal cord itself, as distinct from the nerve roots (radiculopathy) or peripheral nerves (neuropathy). The word says where the problem is, not what caused it.}
}
\newglossaryentry{naturalhistory}{
  name={natural history},
  description={How a disease unfolds over time without treatment. For degenerative cervical myelopathy the natural history is known only roughly, because most modern patients are operated on.}
}
\newglossaryentry{neuroprotection}{
  name={neuroprotection},
  description={Any treatment intended to keep nerve cells alive and functioning under injury, as opposed to relieving the cause of injury. In degenerative cervical myelopathy, neuroprotective drugs are being tested as companions to surgery, not replacements for it.}
}
\newglossaryentry{nonop}{
  name={structured rehabilitation},
  description={Supervised nonoperative care: physiotherapy for strength and balance, activity modification, and sometimes a soft collar, with scheduled reviews to detect deterioration. It is an option only for mild disease and has never been shown to reverse cord damage.}
}
\newglossaryentry{nph}{
  name={normal pressure hydrocephalus},
  description={Enlargement of the brain's fluid spaces in older adults that produces a slow, shuffling, wide-based gait, urinary urgency, and cognitive slowing. Its gait and bladder symptoms overlap with cervical myelopathy, so it is a recognized mimic.}
}
\newglossaryentry{nurick}{
  name={Nurick grade},
  description={A six-step scale (0 to 5) of walking ability in cervical myelopathy, from root symptoms only through difficulty walking, inability to work, needing assistance, to being wheelchair or bed bound. Older and coarser than the mJOA, but simple.}
}
\newglossaryentry{oligodendrocyte}{
  name={oligodendrocyte},
  description={The cell type that makes and maintains myelin in the brain and spinal cord. Their death under chronic compression is a main cause of demyelination, and protecting them is a target of experimental drugs.},
  plural={oligodendrocytes}
}
\newglossaryentry{opll}{
  name={ossification of the posterior longitudinal ligament (OPLL)},
  description={Bone forming inside the ligament that runs down the back of the vertebral bodies, along the front wall of the spinal canal. It narrows the canal from the front, is far more common in East Asian populations, and is one of the causes grouped under degenerative cervical myelopathy.}
}
\newglossaryentry{osteophyte}{
  name={osteophyte},
  description={A bony spur that grows at the edge of a degenerating disc or joint. Combined with a bulging disc it forms the disc-osteophyte complex that indents the front of the spinal cord.},
  plural={osteophytes}
}
\newglossaryentry{paresthesia}{
  name={paresthesia},
  description={Abnormal sensations such as tingling, pins and needles, or burning without an external cause. Hand paresthesias are among the earliest and commonest symptoms of degenerative cervical myelopathy and are easily blamed on carpal tunnel syndrome.},
  plural={paresthesias}
}
\newglossaryentry{prevalence}{
  name={prevalence and incidence},
  description={Prevalence is the share of a population that has a condition at a point in time; incidence is the rate at which new cases arise, often per 100,000 person-years. For degenerative cervical myelopathy both are minimum estimates because most cases are never diagnosed.}
}
\newglossaryentry{pseudarthrosis}{
  name={pseudarthrosis},
  description={Failure of a spinal fusion to heal into solid bone, leaving a persistent false joint that can cause pain or need revision surgery.}
}
\newglossaryentry{qaly}{
  name={quality-adjusted life year (QALY)},
  description={One year of life in perfect health, or the equivalent spread over more years at lower quality. Cost per QALY gained is the standard yardstick for judging whether a treatment is worth its price against a payer's willingness-to-pay threshold.}
}
\newglossaryentry{qmri}{
  name={quantitative MRI},
  description={Advanced MRI measurements that put numbers on cord tissue: cross-sectional area, water diffusion along fibers (diffusion tensor imaging), and myelin content. Under study as a way to detect progression that clinical scales miss.}
}
\newglossaryentry{radiculopathy}{
  name={radiculopathy},
  description={Irritation or compression of a nerve root as it leaves the spinal canal, causing pain, numbness, or weakness in that root's territory in one arm. It is a peripheral nerve problem, not a cord problem, but in someone with cord compression it triples the risk that myelopathy will follow.}
}
\newglossaryentry{recodedcm}{
  name={AO Spine RECODE-DCM},
  description={An international, multi-stakeholder initiative (REsearch Objectives and Common Data Elements for Degenerative Cervical Myelopathy) that agreed a unifying name and definition for the disease, a top-ten list of research priorities, and a minimum data set for studies.}
}
\newglossaryentry{relativerisk}{
  name={relative risk (RR)},
  description={The probability of an outcome in an exposed group divided by the probability in an unexposed group. A relative risk of 3.0 means the outcome is three times as likely; the 95\% confidence interval shows how precisely it is estimated.}
}
\newglossaryentry{reperfusioninjury}{
  name={ischemia-reperfusion injury},
  description={Additional damage that occurs when blood flow suddenly returns to tissue that has adapted to a poor supply, releasing reactive oxygen species. It has been shown in animal models after surgical decompression and may explain why a minority of patients worsen after surgery.}
}
\newglossaryentry{riluzole}{
  name={riluzole},
  description={A drug that dampens excessive glutamate signaling and sodium entry into stressed neurons, approved for ALS. It protected the cord in animal models of myelopathy but did not improve the primary outcome in a large surgical trial.}
}
\newglossaryentry{romberg}{
  name={Romberg test},
  description={The patient stands with feet together and eyes closed; swaying or falling indicates loss of position sense from the legs. It tests the dorsal columns of the spinal cord, which carry position sense and are often damaged in cervical myelopathy.}
}
\newglossaryentry{sci}{
  name={spinal cord injury (SCI)},
  description={Damage to the spinal cord that disrupts movement, sensation, or autonomic function below the injury. Degenerative cervical myelopathy is a slow, nontraumatic form; a compressed cord is also more vulnerable to sudden injury from a fall or whiplash.}
}
\newglossaryentry{spasticity}{
  name={spasticity},
  description={Velocity-dependent stiffness of muscles caused by loss of the brain's inhibitory control over spinal reflexes. In cervical myelopathy it produces stiff, scissoring legs and a slow, unsteady gait.}
}
\newglossaryentry{spinalcanal}{
  name={spinal canal},
  description={The bony tunnel formed by the stacked vertebrae, through which the spinal cord runs bathed in cerebrospinal fluid. Its front wall is the vertebral bodies and discs; its back wall is the laminae and ligamentum flavum.}
}
\newglossaryentry{spondylosis}{
  name={spondylosis},
  description={Age-related wear of the spine: discs lose height and water, bony spurs form at their margins, and joints and ligaments thicken. It is nearly universal after middle age and is the usual cause of the canal narrowing behind degenerative cervical myelopathy.}
}
\newglossaryentry{ssep}{
  name={somatosensory evoked potentials (SSEPs)},
  description={Electrical responses recorded over the spine and scalp after stimulating a nerve in the arm or leg. Slowed conduction shows that the sensory pathways in the cord are damaged, and in people with silent compression it roughly triples the risk of later myelopathy.}
}
\newglossaryentry{stenosis}{
  name={cervical stenosis},
  description={Narrowing of the cervical spinal canal, whether acquired through degeneration or present from birth. A canal narrower than about 13 mm from front to back is usually called congenitally stenotic and leaves little room for later degenerative change.}
}
\newglossaryentry{syringomyelia}{
  name={syringomyelia},
  description={A fluid-filled cavity that forms inside the spinal cord and slowly expands, typically causing loss of pain and temperature sensation in a cape-like pattern over the shoulders and arms. Visible on MRI and readily distinguished from compression.}
}
\newglossaryentry{systematicreview}{
  name={systematic review and meta-analysis},
  description={A study that searches the literature by a preset protocol, appraises every eligible study, and (in a meta-analysis) pools their results into one estimate with a confidence interval. Pooled estimates for degenerative cervical myelopathy often have wide intervals because the underlying studies are small and varied.}
}
\newglossaryentry{t1hypointensity}{
  name={T1 hypointensity},
  description={A dark area inside the cord on T1-weighted MRI. It usually reflects established tissue loss (myelomalacia) and, especially when paired with a T2 bright spot, predicts a poorer response to surgery.}
}
\newglossaryentry{t2hyperintensity}{
  name={T2 hyperintensity},
  description={A bright spot inside the spinal cord on T2-weighted MRI, indicating excess water from swelling, inflammation, or tissue loss. It is common in cervical myelopathy but nonspecific, and on its own it does not settle whether the damage is reversible.}
}
\newglossaryentry{tandemgait}{
  name={tandem gait},
  description={Walking heel to toe along a line. Inability to do it is an early sign of the balance loss in cervical myelopathy and often precedes visible weakness.}
}
\newglossaryentry{tetraplegia}{
  name={tetraplegia},
  description={Loss of movement and sensation in all four limbs from damage to the cervical spinal cord. It is the end point of untreated, progressive cervical myelopathy and the reason the condition is treated as a slow spinal cord injury.}
}
\newglossaryentry{torgpavlov}{
  name={Torg–Pavlov ratio},
  description={The front-to-back diameter of the spinal canal divided by that of the vertebral body at the same level, measured on a side-view X-ray. A ratio below about 0.8 indicates a congenitally narrow canal, independent of X-ray magnification.}
}
\newglossaryentry{tromner}{
  name={Trömner sign},
  description={A cousin of the Hoffmann sign: tapping the pad of the slightly flexed middle finger makes the thumb and index finger flex. Two small studies found it both more sensitive and more specific than the Hoffmann sign.}
}
\newglossaryentry{whitematter}{
  name={white matter},
  description={The outer part of the spinal cord made of long, insulated nerve fibers running up and down. It carries the motor and sensory pathways whose damage produces symptoms below the level of compression.}
}
